For a chain of local subrings ordered by domination, the union is a ring and the union of their maximal ideals is an ideal. Every element outside this ideal is invertible at the stage where it occurs. An element in the ideal cannot become invertible at a later stage, since domination preserves its membership in the maximal ideal. Thus the union is local and dominates the whole chain. Zorn's lemma gives maximal elements, and likewise gives a maximal element above any specified local subring.

Let $A$ be maximal, with maximal ideal $\mathfrak{m}$. Extend $A\to\overline{A/\mathfrak{m}}$ to a maximal pair consisting of a subring of $K$ and a homomorphism to this algebraically closed field. By Lemma 5.19 and Theorem 5.21, the resulting ring is a valuation ring whose maximal ideal is the kernel of the map. It dominates $A$, so maximality makes it equal to $A$.

Conversely, suppose $A$ is a valuation ring and a local subring $B\subseteq K$ dominates it. If $x\in B\setminus A$, then $x^{-1}\in A$ is a nonunit of $A$. Domination makes it a nonunit of $B$, contradicting $x\in B$. Hence $B=A$.
